\section*{NeurIPS Paper Checklist}

\begin{enumerate}

\item {\bf Claims}
    \item[] Question: Do the main claims made in the abstract and introduction accurately reflect the paper's contributions and scope?
    \item[] Answer: \answerYes{}
    \item[] Justification: The abstract distinguishes two types of claims: (1) structural coverage metrics (100\% M-check catch rate, FPR~=~0\% by DFA construction; 100\% CARLog traceability, architectural guarantee) and (2) probabilistic fraud detection metrics (88.0\%$\pm$0.9\% recall, 0.0\%$\pm$0.0\% hard FPR, 59.4\% escalation rate --- multi-seed over 8-attack benchmark, matched exactly in Table~\ref{tab:results}). B7s/B8s per-attack results and the LLM-judge baseline (61\% recall) are seed=42 supplementary results, labeled as such in \S6. The distinction between structural coverage claims and probabilistic detection metrics is elaborated in \S6. The Commerce State Automaton is formally defined (Definition~3.0) with proofs provided for all three theoretical results. No L3 implementation is claimed.

\item {\bf Limitations}
    \item[] Question: Does the paper discuss the limitations of the work performed by the authors?
    \item[] Answer: \answerYes{}
    \item[] Justification: \S\ref{sec:discussion_limits} covers: (1) heuristic uncertainty estimates are boundary-proximity heuristics, not conformal-prediction intervals---the exact conditions required for marginal coverage are stated; (2) 56\% escalation rate (soft-flag FPR; hard FPR is 0.0\%) is a per-persona spend-ceiling overlap limitation, not a calibration artifact; (3) B7s (MCC spoof) at 82.5\% requires merchant-domain knowledge beyond the current allowlist; (4) B4/B5 at 68--70\% require cross-session registries; (5) L3 ceremony tier (ZK proof, selective-disclosure credential) is unimplemented; (6) the commerce automaton is validated against ISO~8583, x402, and Stripe PaymentIntent but not against other payment networks.

\item {\bf Theory assumptions and proofs}
    \item[] Question: For each theoretical result, does the paper provide the full set of assumptions and a complete proof?
    \item[] Answer: \answerYes{}
    \item[] Justification: The main results are stated with explicit assumptions and complete proofs. Theorem~\ref{thm:general_char} (characterization of safe reducers on any finite bounded lattice: safe iff semilattice operation dominating the join; join is the unique minimal safe reducer) has a proof sketch in \S\ref{sec:general_typing} and a full proof in Appendix~\ref{app:general_char}; its assumptions (finite bounded lattice, reducer with $\bot$ as unit, trace semantics S1--S4) are stated in Definitions preceding it. Theorem~\ref{thm:delegation} (delegation closure: composite is fail-closed iff the lattice embedding is top-preserving) is proved in full in \S\ref{sec:general_typing}. Theorems~\ref{thm:type_safety}, \ref{thm:lww_unsound}, \ref{thm:char} (commerce instantiation) are proved in \S\ref{sec:cts}. Proposition~\ref{prop:decidability} (decidability, $O(n)$) and the tamper-evidence and CRDT propositions are proved in the text and appendix. Where a result is a direct specialization (e.g., Theorem~\ref{thm:char} from Theorem~\ref{thm:general_char}), the paper says so explicitly.

\item {\bf Experimental result reproducibility}
    \item[] Question: Does the paper fully disclose all information needed to reproduce the main experimental results?
    \item[] Answer: \answerYes{}
    \item[] Justification: The Reproducibility Statement gives the exact command (\texttt{python -m benchmark.run\_experiment --seed 42 --no-sota}), compute environment (CPU only, $<$30 seconds/seed), and JSON output field names. Markov parameters, amount distributions, perturbation $\sigma$, training seed, threshold formula, and veto conditions are fully specified in the companion AIWILD paper (cited), which contains the benchmark details. All automaton transitions are fully enumerated in Table~1.

\item {\bf Open access to data and code}
    \item[] Question: Does the paper provide open access to data and code?
    \item[] Answer: \answerYes{}
    \item[] Justification: ACP reference implementation and benchmark harness are released as open-source (anonymized URL in supplemental material; full URL in camera-ready). The benchmark is fully deterministic given \texttt{--seed}. The formal automaton definition (Definition~3.0) and CTS channel specification (\S4) are fully self-contained and implementation-independent.

\item {\bf Experimental setting/details}
    \item[] Question: Does the paper specify all training and test details necessary to understand the results?
    \item[] Answer: \answerYes{}
    \item[] Justification: Training data: seed=7, $N=500$ sessions. Test data: seeds 42--46, $N=520$ sessions each. ML threshold: $\mu + 1.1\sigma$ of clean training scores. Veto: score$=0$ and uncertainty$>0.20$. Uncertainty: boundary-proximity heuristic (0.30 if $|$score$-$threshold$|<0.12$, else 0.08). All specified in companion AIWILD paper (\S3 and Appendix~A).

\item {\bf Experiment statistical significance}
    \item[] Question: Does the paper report error bars or appropriate information about statistical significance?
    \item[] Answer: \answerYes{}
    \item[] Justification: All multi-seed results are reported as mean $\pm$ standard deviation across 5 independent seeds (42--46). Standard deviations are population standard deviations, not standard errors---stated in the table caption. The five seeds are fully independent (separate data generation). Rates are bounded well away from 0 and 1 in the reported range, so symmetric intervals are appropriate.

\item {\bf Experiments compute resources}
    \item[] Question: Does the paper provide sufficient information on compute resources needed to reproduce the experiments?
    \item[] Answer: \answerYes{}
    \item[] Justification: CPU only; no GPU required. $<$30 seconds per seed; $<$3 minutes for all five seeds. The Tier~3 LLM Verifier invoked on $<$3\% of actions in the full pipeline is bypassed by \texttt{--no-sota} in all reported experiments. Stated in the Reproducibility Statement.

\item {\bf Code of ethics}
    \item[] Question: Does the research conform with the NeurIPS Code of Ethics?
    \item[] Answer: \answerYes{}
    \item[] Justification: This work provides formal safety guarantees and defensive tooling for financial AI systems. All benchmark data is synthetic. No production payment credentials or working exploit code are released. The attack taxonomy is presented at the conceptual level required to motivate the defensive framework.

\item {\bf Broader impacts}
    \item[] Question: Does the paper discuss both potential positive and negative societal impacts?
    \item[] Answer: \answerYes{}
    \item[] Justification: See the Broader Impact section following the Conclusion. Positive: formal safety guarantees at the protocol layer (L1) reduce the window for financial fraud in agentic systems; the decidability result enables static compliance tooling. Negative: publication of the automaton model could inform adversarial agents that probe state transition boundaries; mitigated by the verifier being released alongside the taxonomy.

\item {\bf Safeguards}
    \item[] Question: Does the paper describe safeguards for responsible release of data or models with high risk for misuse?
    \item[] Answer: \answerYes{}
    \item[] Justification: The formal framework defines safety properties; the release includes the verifier (defensive tool) and the synthetic generator, not production payment attack implementations. No live API credentials, no production merchant exploit code, and no B5 settled-key replay implementations are released.

\item {\bf Licenses for existing assets}
    \item[] Question: Are creators of assets used in the paper properly credited and license terms mentioned?
    \item[] Answer: \answerYes{}
    \item[] Justification: LangGraph~\cite{langgraph} (Apache 2.0); ISO~8583 is a published standard (cited); x402~\cite{x402} (open protocol, arXiv); Stripe~\cite{stripe} (public API documentation). All appear in the References section.

\item {\bf New assets}
    \item[] Question: Are new assets introduced in the paper well documented?
    \item[] Answer: \answerYes{}
    \item[] Justification: The Commerce State Automaton (Definition~3.0), CTS channel type (Definition~4.1), and CARLog (Definition~4.2) are formally defined with full enumeration of states, transitions, and operators. The ACP reference implementation is documented in \S4 and released open-source under MIT license.

\item {\bf Crowdsourcing and research with human subjects}
    \item[] Question: Does the paper include full details of crowdsourcing or human subjects research?
    \item[] Answer: \answerNA{}
    \item[] Justification: All experimental data is synthetically generated. No human subjects, no crowdsourcing, no real user data.

\item {\bf Institutional review board (IRB) approvals}
    \item[] Question: Does the paper describe potential risks to study participants and IRB approvals?
    \item[] Answer: \answerNA{}
    \item[] Justification: No human subjects research. Not applicable.

\item {\bf Declaration of LLM usage}
    \item[] Question: Does the paper describe the usage of LLMs if it is an important, original, or non-standard component of the core methods?
    \item[] Answer: \answerYes{}
    \item[] Justification: Two LLM uses are disclosed. (1) The Tier~3 LLM Verifier (\S4) uses \texttt{claude-haiku-4-5-20251001} for high-uncertainty actions ($<3\%$ of actions), described in \S4 (``Ceremony Tiers and Escalation Protocol''). (2) The full-session LLM-as-judge baseline (\S6, Table~\ref{tab:results}) uses \texttt{claude-sonnet-4-6} via AWS Bedrock to establish the L0 detection ceiling (61\% recall, seed=42)---an evaluation component, not a method component. LLMs are \emph{not} used in the formal automaton, the CTS channel type, or the CARLog---those are provably correct by construction. LLMs were also used for writing assistance in preparing this manuscript, which does not require declaration per the NeurIPS handbook.

\end{enumerate}
